\documentclass[12pt]{amsart}
\usepackage[margin=1in]{geometry}
\usepackage{amsmath,amssymb,amsthm}
\usepackage{mathtools}
\usepackage{hyperref}

\theoremstyle{plain}
\newtheorem{theorem}{Theorem}[section]
\newtheorem{proposition}[theorem]{Proposition}
\newtheorem{corollary}[theorem]{Corollary}
\newtheorem{lemma}[theorem]{Lemma}
\newtheorem{definition}[theorem]{Definition}
\theoremstyle{remark}
\newtheorem*{remark}{Remark}

\title[Dissipation Succeeds Where Unitarity Cannot]{Dissipation Succeeds Where Unitarity Cannot:\\
A Mean-Field Realization of $B_X$}
\author{Jeremy Ebert}
\address{Franklin County, Ohio, USA}
\date{2026}

\begin{document}
\maketitle

\begin{abstract}
\cite{Ebert2026FourModes} proves that no \emph{unitary} generator --- quadratic, quartic-diagonal (Kerr), or any natural quaternion-motivated combination of the shape modes with a second oscillator pair --- realizes $B_X$'s classical flow, $x(\phi)=x_0e^\phi$, $y(\phi)=y_0e^{-\phi}$, $Y=\sqrt{xy}$ exactly constant. That result concerns closed, reversible (Hamiltonian-generated) dynamics only. This note shows that a genuinely \emph{dissipative} process --- outside the class \cite{Ebert2026FourModes} ruled out --- succeeds, exactly, at the level of mean photon numbers: coupling mode $1$ to a gain reservoir and mode $2$ to a loss reservoir, at equal rates, via the standard Lindblad dissipators $\mathcal D[a_1^\dagger]$, $\mathcal D[a_2]$, gives $\langle n_1(\phi)\rangle=(n_1(0)+1)e^\phi-1$, $\langle n_2(\phi)\rangle=n_2(0)e^{-\phi}$ --- $B_X$'s flow exactly, not approximately --- with $(\langle n_1\rangle+1)\langle n_2\rangle$ exactly constant. We verify this both from the closed moment equations and from direct integration of the full master equation (not assuming factorization), confirming the two modes remain in a product state throughout. We are explicit and emphatic about what this does not establish: Lindblad evolution is irreversible and non-unitary, degrades pure states into mixed ones, and is a completely different mathematical object (a superoperator on density matrices) from a Lie-algebra generator whose exponential would give a genuine quantum representation of the classical symmetry. This result does not contradict \cite{Ebert2026FourModes}; it answers a different, complementary question --- is $B_X$'s classical flow realizable by \emph{any} quantum process --- in the affirmative, while leaving the original question --- is there a unitary quantum $B_X$ --- exactly as open as \cite{Ebert2026FourModes} left it.
\end{abstract}

\section{What the Unitary No-Go Left Unexplored}

\cite{Ebert2026FourModes} rules out closed-system, Hamiltonian-generated dynamics: single or independent-sum passive/active bilinears, and Kerr-type diagonal additions to them. All of these are \emph{unitary} --- generated by $\exp(-iHt)$ or $\exp(rK)$ for Hermitian or anti-Hermitian $H,K$. This leaves an entire further class of Gaussian-preserving quantum operations untouched: dissipative channels, described not by a Hamiltonian but by a Lindblad master equation, which are not reversible and do not preserve the purity of a state. We test whether this broader class succeeds where the narrower one provably cannot.

\section{Engineered Gain and Loss}

\begin{definition}
Let $\rho$ evolve by $\dot\rho=\kappa\mathcal D[a_1^\dagger]\rho+\kappa\mathcal D[a_2]\rho$, where $\mathcal D[L]\rho:=L\rho L^\dagger-\tfrac12\{L^\dagger L,\rho\}$: mode $1$ coupled to a gain reservoir, mode $2$ to a loss reservoir, at equal rates $\kappa$.
\end{definition}

\begin{lemma}[Exact linear moment equations]\label{lem:moments}
\begin{equation}
\frac{d\langle n_1\rangle}{dt}=\kappa(\langle n_1\rangle+1), \qquad \frac{d\langle n_2\rangle}{dt}=-\kappa\langle n_2\rangle. \tag{1}
\end{equation}
\end{lemma}

\begin{proof}
For a single Lindblad term with jump operator $L$, $\frac{d\langle O\rangle}{dt}=\langle L^\dagger OL\rangle-\tfrac12\langle L^\dagger LO\rangle-\tfrac12\langle OL^\dagger L\rangle$. For $L=\sqrt\kappa\,a^\dagger$, $O=n=a^\dagger a$: using $na^\dagger=a^\dagger(n+1)$ and $aa^\dagger=n+1$, $L^\dagger OL=\kappa\,a\,n\,a^\dagger=\kappa(n+1)^2$; $L^\dagger LO=OL^\dagger L=\kappa\,a a^\dagger n=\kappa n(n+1)$. So $\frac{d\langle n\rangle}{dt}=\kappa\big[(n+1)^2-n(n+1)\big]=\kappa(n+1)$, giving the first equation of (1). For $L=\sqrt\kappa\,a$: using $na=a(n-1)$, $L^\dagger OL=\kappa\,a^\dagger n a=\kappa\,n(n-1)$; $L^\dagger LO=OL^\dagger L=\kappa\,a^\dagger a\,n=\kappa n^2$. So $\frac{d\langle n\rangle}{dt}=\kappa[n(n-1)-n^2]=-\kappa n$, giving the second equation.
\end{proof}

\begin{theorem}[Exact agreement with $B_X$]\label{thm:agree}
The solution of (1), $\langle n_1(t)\rangle=(n_1(0)+1)e^{\kappa t}-1$, $\langle n_2(t)\rangle=n_2(0)e^{-\kappa t}$, matches $B_X$'s classical flow $x(\phi)=x_0e^\phi$, $y(\phi)=y_0e^{-\phi}$ exactly under $x=n_1+1$, $y=n_2$, $\phi=\kappa t$, and
\begin{equation}
\big(\langle n_1(t)\rangle+1\big)\langle n_2(t)\rangle = (n_1(0)+1)n_2(0) \tag{2}
\end{equation}
is exactly constant.
\end{theorem}

\begin{proof}
Direct substitution: $(\langle n_1\rangle+1)\langle n_2\rangle=(n_1(0)+1)e^{\kappa t}\cdot n_2(0)e^{-\kappa t}=(n_1(0)+1)n_2(0)$, the exponentials cancelling exactly.
\end{proof}

\subsection*{Numerical Verification}
With $n_1(0)=3$, $n_2(0)=2$ (so $Y^2=(n_1(0)+1)n_2(0)=8$), (1)--(2) give $Y^2=8.0000$ at $\phi=0,0.5,1,1.5,2$ to displayed precision, confirmed by direct evaluation of the closed forms. We then integrated the \emph{full} master equation (not the mean-value ODEs) on a truncated ($N=15$ per mode) two-mode Fock space from the product Fock state $|3,2\rangle$: $\langle(n_1+1)n_2\rangle$, computed from the full density matrix, matches $(\langle n_1\rangle+1)\langle n_2\rangle$ (the factorized product of separately-computed means) exactly at every sampled $t\in\{0,0.15,0.3,0.45\}$, and $\operatorname{tr}\rho=1$ to $10^{-6}$ throughout. A small drift in $Y^2$ (from $8.0000$ to $7.9984$ at the latest time sampled) is attributable to the finite $N=15$ truncation becoming inadequate as $\langle n_1\rangle$ grows past $5$, not to any breakdown of (2) itself.

\begin{remark}[No entanglement develops]
The exact match between the full-$\rho$ moment $\langle(n_1+1)n_2\rangle$ and the factorized $(\langle n_1\rangle+1)\langle n_2\rangle$ confirms the two modes remain in a product state throughout: gain and loss act independently on their own modes, with no interaction term to generate correlations between them.
\end{remark}

\section{What This Does Not Establish}

We state this plainly rather than let the clean agreement above imply more than it does.

\begin{itemize}
\item[(i)] \textbf{This is not a unitary generator.} A Lindblad master equation is a superoperator acting on density matrices, not an element of $\mathfrak{so}(2,1)$ or any Lie algebra whose exponential gives a reversible transformation. There is no single operator $\hat B_X$ here whose flow this is; there is a dissipator, and dissipators are not invertible.
\item[(ii)] \textbf{Purity is not preserved.} Starting from the pure state $|3,2\rangle$, the evolution described here generically produces a mixed state (the gain and loss reservoirs inject genuine quantum noise). $B_X$'s classical flow is a well-defined trajectory on the cone; this construction reproduces it only in the sense of first moments of a decohering quantum state, not as a coherent quantum trajectory.
\item[(iii)] \textbf{This does not contradict \cite{Ebert2026FourModes}.} That note's Lemmas and Proposition concern generators of unitary evolution exclusively. Nothing here is a counterexample to any claim made there; it is a genuinely different class of quantum operation, outside the scope of what was ruled out.
\end{itemize}

\begin{remark}[What would be needed to go further]
Whether \emph{some} unitary construction --- necessarily non-Gaussian and not reducible to Kerr-enhanced Gaussian coupling, per \cite{Ebert2026FourModes} --- could reproduce $B_X$'s flow coherently remains exactly as open as it was. What this note adds is a existence result of a different kind: the flow itself is physically achievable by a quantum process, which was not previously established in any form, coherent or not.
\end{remark}

\section{Discussion}

This result reuses, in a new role, exactly the Lindblad machinery already developed for the boson-bath calculations of \cite{Ebert2026Bath}: there, gain and loss dissipators were used to explain \emph{why} a coordinate relaxes toward an equilibrium set by a bath. Here, the same machinery is repurposed constructively: rather than asking what a bath does to an existing coordinate, we ask whether an engineered bath can be built to reproduce a specific target flow, and for $B_X$ specifically, it can, exactly, at the mean-field level. The honest reading is that this substantially reframes rather than resolves the original question of \cite{Ebert2026Quantum}: the search for a coherent, unitary quantum $B_X$ remains unsuccessful, but the question of whether $B_X$'s classical content has \emph{any} quantum-mechanical existence now has a clear, verified, affirmative answer, of a kind distinct from what a symmetry generator would provide.

\section*{AI Assistance Disclosure}

Large language model tools (Anthropic's Claude) were used in the derivation, symbolic and numerical cross-checking, and \LaTeX{} preparation of this note. This includes: the direct operator-algebra derivation and verification of Lemma~2.2's two moment equations; the closed-form verification of Theorem~2.3 and the exact cancellation in (2); the full master-equation simulation verifying the closed-form solution against direct numerical integration of the density matrix (not merely the mean-value ODEs), including the factorization check ruling out spurious inter-mode correlations; and identification of the finite-truncation origin of the small late-time drift in the numerical $Y^2$. All mathematical claims were independently verified by the author.

\begin{thebibliography}{9}
\bibitem{Ebert2026Cone} J.~Ebert, \emph{The Dynamic Slicing Operator on the AM--GM Cone: A Quadric Geometric Framework for Keplerian Orbits and Perturbed Trajectories}, Preprint (2026).
\bibitem{Ebert2026Bath} J.~Ebert, \emph{The Cone in a Bath: Boson, Gravitational, and Plasma Relaxation of the Semi-Major Axis}, Preprint (2026).
\bibitem{Ebert2026Quantum} J.~Ebert, \emph{Two Oscillators, Two Algebras: SU(2) and SU(1,1) Realizations of the Cone's Symmetry, and the Fate of the Power Map}, Preprint (2026).
\bibitem{Ebert2026FourModes} J.~Ebert, \emph{Four Modes, Still No $B_X$: Passive Correlation, Active Growth, and Why Kerr Terms Cannot Bridge Them}, Preprint (2026).
\end{thebibliography}

\end{document}
